\documentclass[11pt]{article}
\usepackage[margin=1in]{geometry}
\usepackage{amsmath,amssymb,booktabs}
\usepackage[T1]{fontenc}
\usepackage{xurl}
\usepackage[hidelinks]{hyperref}
\begin{document}
\title{Historical joint calibration draft: not the selected Attempt 209 paper}
\author{}\date{}\maketitle
\noindent\textbf{Historical draft.} This text describes a later joint fit, not the fixed-hardening Attempt 209 calibration selected for publication. Do not pair its parameters or optimization narrative with the selected Attempt 209 figure bundle. A matching manuscript update is still required.\par
\noindent\textit{Draft for transcription. Equation numbers and citations should be adapted to the manuscript. Runs 1 and 3 enter the joint fit; Run 2 is an exploratory frozen-parameter check, not a pristine independent test. Editorial qualifications are collected at the end.}

\section{Model calibration}
\subsection{Parameters and deformation history}
The calibration treats growth inputs as effective parameters of the shell model, not as independently measured material constants. For each layer $\ell\in\{t,b\}$, an amplitude $g_\ell$ and an orthotropy coefficient $o_\ell$ define the two local directional growth inputs,
\begin{equation}
 g_{\ell,1}=g_\ell(1+o_\ell),\qquad
 g_{\ell,2}=g_\ell(1-o_\ell).
 \label{eq:directional}
\end{equation}
The inputs are modulated by the prescribed track profile and by a face-shared plastic-hit history. For $n$ previous counted hits, the history multiplier is
\begin{equation}
 q(n)=f+(1-f)\exp(-\beta n),
 \label{eq:hardening}
\end{equation}
where $\beta$ is the decay rate and $f$ is the asymptotic floor. The multiplier acts on the growth increment for each hit; it is not a post hoc scaling of the final deformation.

The broader search allowed the six parameters $(g_t,g_b,o_t,o_b,\beta,f)$ to vary. In the final reduced search, $g_b$ was fixed at zero and $o_b$ was set to zero because it has no effect when $g_b=0$. The four active parameters and their final search bounds are shown in Table~\ref{tab:parameters}. The two positive parameters $g_t$ and $\beta$ were searched logarithmically, and $o_t$ and $f$ linearly. Fixing $g_b=0$ is a model-reduction choice within the explored search, not evidence that physical bottom-layer growth is absent.

\begin{table}[ht]
\centering
\caption{Selected parameters and bounds in the final reduced search. All values are dimensionless.}
\label{tab:parameters}
\begin{tabular}{llll}\toprule
Parameter & Lower bound & Upper bound & Selected value \\\midrule
$g_t$ & $5\times10^{-5}$ & $10^{-2}$ & $1.45736\times10^{-4}$ \\
$o_t$ & $-1$ & $1$ & $-0.47931$ \\
$\beta$ & $0.01$ & $1.0$ & $0.49094$ \\
$f$ & $0$ & $0.5$ & $0.18692$ \\\midrule
$g_b$ & \multicolumn{2}{l}{fixed} & $0$ \\
$o_b$ & \multicolumn{2}{l}{inactive when $g_b=0$} & $0$ \\\bottomrule
\end{tabular}
\end{table}

The cases use a $200\times300\times0.77$ mm sheet, a $2$ mm effective track width, and a mesh with 15,251 vertices and 30,000 triangular elements. The recorded numerical material settings use $E=1$ and $\nu=0.33$; the former is not a measured modulus. The prescribed deformation history is accumulated before a final unconstrained equilibrium solve. The equilibrium procedure uses an HLBFGS-based warm-up followed by sparse Newton steps with adaptive cutback, with a gradient-norm tolerance of $5\times10^{-14}$ in internal solver units. A common parameter set is fitted to Runs 1 and 3, which differ in prescribed path and cycle count.

\subsection{Geometry comparison and objective}
The measured top surface is approximated as the shell midsurface. The simulation is brought into a material-reference frame using matched deformed and undeformed vertices. The scan and simulation are then leveled and centered against a flat shared-footprint reference without rigidly fitting the scan to the prediction. Public axes are $x$ for sheet width, $z$ for length and $y$ for height.

For case $c$, let $\widehat{\mathbf p}_{ci}$, $i=1,\ldots,N_c$, be the retained aligned scan points, $\widehat S_c(\boldsymbol\theta)$ the framed simulated triangle surface, and $\widehat T_c$ the triangle surface formed from the scan footprint. Let $\widehat{\mathbf q}_{ci}(\boldsymbol\theta)$ be a point on the simulated surface sampled at the corresponding footprint location. The directed distances and the case discrepancy are
\begin{align}
 d_{ci}(\boldsymbol\theta)&=\min_{\mathbf q\in\widehat S_c(\boldsymbol\theta)}
          \lVert\widehat{\mathbf p}_{ci}-\mathbf q\rVert_2,\\
 r_{ci}(\boldsymbol\theta)&=\min_{\mathbf p\in\widehat T_c}
          \lVert\widehat{\mathbf q}_{ci}(\boldsymbol\theta)-\mathbf p\rVert_2,\\
 D_c(\boldsymbol\theta)&=\left\{\frac12\left[
 \frac{1}{N_c}\sum_{i=1}^{N_c}d_{ci}^{2}
 +\sum_{i=1}^{N_c}a_{ci}r_{ci}^{2}\right]\right\}^{1/2}.
 \label{eq:case-objective}
\end{align}
Here $a_{ci}$ are normalized barycentric dual-area weights computed from the projected scan triangles. The scan triangulation excludes long edges that could bridge gaps. The predicted points use the actual model connectivity; where vertical interpolation on the deformed projection is unavailable, material-coordinate interpolation is used. Thus both directions sample the shared scan footprint, although the forward nearest-point query is made against the framed model surface. Distances are in millimetres. The use of triangle-surface distances avoids the nonzero sampling and boundary floor of the earlier point-cloud reverse comparison.

Runs 1 and 3 have equal case weights. The joint objective minimized in the final search is
\begin{equation}
 J(\boldsymbol\theta)=\left[\frac{D_1(\boldsymbol\theta)^2+D_3(\boldsymbol\theta)^2}{2}\right]^{1/2}.
 \label{eq:joint-objective}
\end{equation}
Run 2 does not enter Eq.~\eqref{eq:joint-objective}. This equal-weight joint objective replaces the earlier single-case pilot objective.

\subsection{Bayesian search}
An earlier joint search was followed by two separate comparison arms with the same initial design: a success-aware search and a failure-imputation search. A subsequent bottom-growth extension inherited the success-aware arm's records. Some earlier evaluations were later rescored using the revised objective, so their displayed scores did not necessarily guide their original proposals. These searches used different parameter bounds or policies and are not one uninterrupted trajectory.

The final search was warm-started using ten eligible valid results with zero bottom growth from the preceding extension. Its first four candidates were prescribed controls obtained by setting bottom growth to zero in selected earlier candidates and re-evaluating them. Subsequent candidates were chosen in batches of four by a success-aware acquisition: a Mat\'ern-$5/2$ Gaussian-process surrogate was fitted to $-\log[J/(1\,\mathrm{mm})]$ for valid evaluations, and expected improvement was weighted by a separate estimated probability of solver success. Failed evaluations informed the success model but were not assigned artificial objective values in the objective surrogate. Candidates in a batch were selected greedily using predicted-mean fantasies, subject to a minimum separation in normalized parameter space.

The final stage had a maximum allocation of 96 new evaluations. It stopped under the planned improvement rule after 12 waves (48 evaluations): 44 were valid, three failed in the numerical solver and one timed out. The selected candidate had the smallest valid joint objective, $J=0.488$ mm. Stopping by the improvement rule does not establish a global optimum or parameter convergence. The previous, broader six-parameter search and the final four-parameter continuation are distinct stages and should not be reported as a single uninterrupted four-parameter search.

\subsection{Fit and parameter assessment}
Table~\ref{tab:fits} compares the selected fit with the previously selected parameter set, evaluated using the same revised surface metric. On the two fitted cases, the joint RMS decreased from $0.674$ to $0.488$ mm. Most of this reduction is associated with Run 3: its case RMS decreased from $0.719$ to $0.290$ mm, whereas the Run 1 case RMS remained approximately $0.626$ mm. These are geometric discrepancies for one scan per case, not uncertainty estimates from independent replicates.

\begin{table}[ht]
\centering
\caption{Case-wise symmetric surface RMS and scan-to-model 95th-percentile distance ($P_{95}$), in mm. The previous and selected values use the revised metric.}
\label{tab:fits}
\begin{tabular}{llrrrr}\toprule
Case & Role & Previous RMS & Selected RMS & Previous $P_{95}$ & Selected $P_{95}$ \\\midrule
Run 1 & joint fit & 0.626 & 0.627 & 1.185 & 1.177 \\
Run 2 & exploratory check & 0.678 & 0.641 & 1.324 & 1.275 \\
Run 3 & joint fit & 0.719 & 0.290 & 1.389 & 0.563 \\\bottomrule
\end{tabular}
\end{table}

The four selected active parameters lie inside their final search bounds; $g_b=0$ was fixed at the boundary. Other combinations give similar but not identical fit quality. For example, one valid alternative with $(\beta,f)=(0.157,0.449)$ has $J=0.508$ mm, compared with $0.488$ mm at $(0.491,0.187)$; its growth amplitude and orthotropy also differ. This is consistent with parameter compensation, but it does not identify a unique set of material parameters. Compared with an earlier final-stage control at $(\beta,f)=(0.300,0.128)$, the selected candidate reduces $J$ from $0.507$ to $0.488$ mm and the Run 1 RMS from $0.673$ to $0.627$ mm, while increasing the Run 3 RMS from $0.246$ to $0.290$ mm. Several parameters changed together, so this improvement cannot be attributed to freeing the hardening parameters alone. Post-hoc GP response slices describe sensitivity conditional on the other parameters and the fitted surrogate; their predictive bands are not confidence intervals for physical parameters.

Run 2 was evaluated using the frozen selected parameters and was not part of the joint objective. Its symmetric RMS was $0.641$ mm, compared with $0.678$ mm for the previous parameter set under the same revised metric. However, earlier Run 2 results were already known, the case was inspected during the search, and an initial time-limited solve was followed by a researcher-approved longer-limit retry. This comparison is therefore exploratory, not a blinded or one-shot validation. The old $1.250$ mm Run 2 value used a superseded metric and is not directly comparable.

Full-field height maps, scan-to-model distance distributions and fixed-$z$ sections complement the objective. The signed height residual $\Delta y(x,z)=y_{\rm model}(x,z)-y_{\rm experiment}(x,z)$ is not the Euclidean distance used in Eq.~\eqref{eq:case-objective}. ``Analytical'' in the figures denotes the numerical shell-model prediction, not a closed-form solution.

\section*{Editorial notes and evidence (not manuscript text)}
\begin{itemize}
\item Public labels are Run 1 = historical \texttt{run6}, Run 2 = historical \texttt{run9}, Run 3 = historical \texttt{run10}. Preserve the historical identifiers in source data and records.
\item The figure drafts and captions are in \path{calibration/paper_figures/current_joint_calibration/}; \path{figure_data_and_provenance.json} in that directory records input hashes, plotted values and limits. Figures use saved evaluations only.
\item The post-hoc GP uses 54 unique valid points; training through wave 8 and checking 16 later valid wave 9--12 points gives log-objective RMSE $0.231$ and 15/16 nominal 95\% interval coverage. The later points are adaptive and not an independent experimental test.
\item The physical gripped edge and worked-face interpretation remain unconfirmed. Verify laboratory fixture records and scan-to-midsurface approximation before making physical claims. The mesh size is not a mesh-convergence study.
\item Primary records: \path{docs/calibration/bayesian_calibration/convention_corrected_investigation.md}, \path{calibration/bo/stage4_reduced4.json}, \path{run/joint_stage4_reduced4/summary.json}, and \path{run/run9_interim/joint_stage4_reduced4_a0016_long/evaluations/attempt_0001/result.json}. The original pilot draft is preserved separately as \path{calibration_manuscript_draft_legacy_pilot.tex}.
\item Add appropriate bibliography entries for the optimization and numerical solver before publication. The figure captions contain selection, interpolation, display-saturation and dependence disclosures.
\end{itemize}
\end{document}
